% Propietario reproducido por unidad demostrativa; véase PROPIETARIOS_CONTINUO_UNIVERSO.json.
% Origen: XI ES CUMPLIMIENTO_EDITORIAL_20260928, sections/supervivencia_punto_fijo.tex:1--146.
% Cuerpos y pruebas conservados; sólo namespace, rutas y remisiones contextuales declaradas.
% BEGIN REV09_SUPERVIVENCIA_PUNTO_FIJO
\subsection{Supervivencia por poda de prolongaciones}
\label{hmtfg:base:corpus9:xi:supervivencia}

La memoria retrospectiva distingue las historias registradas por las
operaciones APP--TRIT--TPK; la compatibilidad prospectiva organiza sus
prolongaciones admisibles. Sea \(S\) el espacio de estados completos
alcanzables mediante esas operaciones. Cada estado conserva hojas, residuos,
cocientes, orientación, acarreo, memoria y datos de frontera; cuando la
transición depende de la profundidad, ésta forma parte del estado.
Sea \(E\subseteq S\times S\) la relación efectiva de prolongación y
\(A\subseteq S\) el conjunto admisible determinado por la compatibilidad de
esos registros. Así, la admisibilidad se decide mediante las operaciones y
las condiciones de frontera que generan la historia.

Para \(X\subseteq S\), definimos
\begin{equation}
 \operatorname{Pre}_{E}(X)
 =\{s\in S:\exists t\in X,\ (s,t)\in E\},
 \qquad
 F_{\mathrm{surv}}(X)=A\cap\operatorname{Pre}_{E}(X).
 \label{hmtfg:base:corpus9:xi:operador-poda}
\end{equation}
La sucesión de poda se obtiene por
\begin{equation}
 V_0=A,\qquad
 V_{n+1}=F_{\mathrm{surv}}(V_n),\qquad
 V_\infty=\bigcap_{n\geq0}V_n.
 \label{hmtfg:base:corpus9:xi:sucesion-poda}
\end{equation}

\begin{proposition}[Supervivencia y mayor punto fijo]
\label{hmtfg:base:corpus9:xi:mayor-punto-fijo}
El conjunto \(V_n\) consta exactamente de los estados de \(A\) que admiten
una prolongación de \(n\) pasos cuyos estados pertenecen a \(A\). La
sucesión \((V_n)\) es decreciente. Si cada estado tiene un número finito
de sucesores admisibles, entonces
\[
 V_\infty=F_{\mathrm{surv}}(V_\infty)
          =\nu F_{\mathrm{surv}},
\]
donde \(\nu F_{\mathrm{surv}}\) denota el mayor punto fijo por inclusión.
Sus elementos son precisamente los estados que admiten una prolongación
infinita compatible.
\end{proposition}

\begin{proof}
Para \(n=0\), una prolongación de longitud cero consiste en el propio
estado de \(A\). Si la caracterización vale para \(n\), un estado pertenece
a \(V_{n+1}\) exactamente cuando pertenece a \(A\) y tiene un sucesor en
\(V_n\); anteponer esa arista produce una prolongación admisible de
\(n+1\) pasos. Esta equivalencia prueba la afirmación por inducción.
Además, \(V_1\subseteq V_0\), y la monotonía de
\(F_{\mathrm{surv}}\) propaga \(V_{n+1}\subseteq V_n\).

Sea \(s\in V_\infty\), y escribamos
\(E(s)=\{t\in S:(s,t)\in E\}\). Como \(s\in V_{n+1}\) para todo \(n\),
los conjuntos \(E(s)\cap V_n\) son no vacíos, finitos y decrecientes.
Su intersección contiene un elemento \(t\), que pertenece a
\(V_\infty\) y satisface \((s,t)\in E\). Por tanto,
\(V_\infty\subseteq F_{\mathrm{surv}}(V_\infty)\).
Recíprocamente, la monotonía da
\(F_{\mathrm{surv}}(V_\infty)\subseteq V_{n+1}\) para todo \(n\);
la imagen también está contenida en \(A=V_0\), de modo que pertenece a
\(V_\infty\). Esto prueba la igualdad de punto fijo.

Si \(W=F_{\mathrm{surv}}(W)\), entonces \(W\subseteq A=V_0\).
De \(W\subseteq V_n\) se deduce
\(W=F_{\mathrm{surv}}(W)\subseteq F_{\mathrm{surv}}(V_n)=V_{n+1}\).
La inducción implica \(W\subseteq V_\infty\), que prueba la maximalidad.
Finalmente, el árbol de prolongaciones de un estado de \(V_\infty\)
es finitamente ramificado y tiene un nivel no vacío a toda profundidad.
El lema de König proporciona una rama infinita. En sentido inverso, los
prefijos de cualquier prolongación infinita sitúan su estado inicial
en todos los \(V_n\).
\end{proof}

\paragraph{Ramificación y estabilización.}
La hipótesis de finitud tiene un control explícito. Considérese una raíz
con una rama terminal distinta de cada longitud entera positiva, con todas
las ramas admisibles y disjuntas después de la raíz. La raíz admite
prolongaciones de cualquier longitud finita y pertenece a todos los
\(V_n\); cada sucesor suyo pertenece a una rama de longitud acotada.
Por ello la raíz carece de prolongación infinita y de sucesor en
\(V_\infty\). Este árbol tiene ramificación infinita en la raíz.
En cambio, si \(A\) es finito, cada inclusión estricta
\(V_{n+1}\subsetneq V_n\) elimina al menos un estado. La poda alcanza
un punto fijo después de, como máximo, \(|A|\) eliminaciones estrictas
de capas. Esta cota pertenece al dominio finito \(A\).

\paragraph{Unicidad de la continuación.}
Para \(s\in V_\infty\), los sucesores que conservan supervivencia forman
\[
 E(s)\cap V_\infty.
\]
La selección del próximo estado es única exactamente cuando
\begin{equation}
 \#\bigl(E(s)\cap V_\infty\bigr)=1.
 \label{hmtfg:base:corpus9:xi:unicidad-sucesor}
\end{equation}
Una prolongación infinita es única cuando esta condición se mantiene
en cada estado de la ruta. En efecto, cada sucesor superviviente tiene
una prolongación infinita por la proposición anterior; dos sucesores
distintos producen dos rutas distintas. Cuando el sucesor es único en
cada paso, los prefijos quedan determinados sucesivamente.
La pluralidad de sucesores viables expresa multiplicidad de
prolongaciones dentro del mismo dominio de supervivencia.

\begin{proposition}[Conjugación de las dos direcciones de prolongación]
\label{hmtfg:base:corpus9:xi:conjugacion-supervivencia}
Sea \(C:S\to S\) una involución tal que \(C(A)=A\) y
\begin{equation}
 (x,y)\in E\quad\Longleftrightarrow\quad(Cy,Cx)\in E.
 \label{hmtfg:base:corpus9:xi:involucion-prolongaciones}
\end{equation}
Denotemos por \(V_n^+\) la poda para \(E\) y por \(V_n^-\) la poda
para \(E^{-1}\), ambas con conjunto inicial \(A\). Entonces
\[
 C(V_n^+)=V_n^-\quad(n\geq0),\qquad
 C(V_\infty^+)=V_\infty^-.
\]
\end{proposition}

\begin{proof}
Si \(s_0,\ldots,s_n\) es una ruta admisible para \(E\), la sucesión
\(Cs_0,\ldots,Cs_n\) es una ruta admisible para \(E^{-1}\):
para cada \(j\), la hipótesis convierte
\((s_j,s_{j+1})\in E\) en \((Cs_{j+1},Cs_j)\in E\).
Todos sus estados pertenecen a \(A\). La misma operación aplicada otra
vez recupera la ruta inicial porque \(C^2=I\); por tanto, establece una
bijección entre las prolongaciones finitas de ambas direcciones.
La caracterización de \(V_n^\pm\) por rutas da
\(C(V_n^+)=V_n^-\). Al ser \(C\) biyectiva, conmuta con la
intersección de estos conjuntos, y se obtiene la igualdad de
\(V_\infty^\pm\).
\end{proof}

La involución transporta el conjunto superviviente de una dirección al
de la dirección conjugada. La prolongabilidad infinita en ambas
direcciones corresponde a \(V_\infty^+\cap V_\infty^-\); la unicidad
de una ruta conserva el criterio de
\eqref{hmtfg:base:corpus9:xi:unicidad-sucesor}. La retención de una genealogía
y la restricción de su árbol de continuaciones son dos operaciones
tipadas sobre los mismos registros. La asignación de una entropía a
esas lecturas requiere, además, la medida de historias correspondiente.
% END REV09_SUPERVIVENCIA_PUNTO_FIJO
